% NOTE: This file is not included by main.tex. The active source is chapters_new.tex and chapters_6_9.tex.
% ==================== 第7章 ====================
\chapter{实验结果}
\label{ch:results}

\section{确定分支：DocFinQA}

\subsection{主结果}

表\ref{tab:docfinqa}报告DocFinQA的300题固定样本。最终确定分支使用16段预算，并在阅读器不充分或置信度低于0.8时回退到固定16段证据集。在DeepSeek上，该分支取得0.457精确匹配，平均使用10.69段；固定16段阅读器取得0.390，平均16段；固定8段阅读器取得0.367，平均8段。无回退的预算8规划器取得0.417，平均4.56段；预算12取得0.427，平均5.82段；预算16取得0.460，平均6.85段。回退在24\%的问题上使用更大证据集，弥补了规划器单独的准确率缺口。

在GLM上，该分支取得0.497，平均10.31段；固定16段阅读器取得0.480，平均16段；固定8段阅读器取得0.440，平均8段。无回退的预算8规划器取得0.450，平均4.15段；预算12取得0.473，平均5.11段；预算16取得0.470，平均5.88段。回退分支比固定16段阅读器高0.017，平均少用5.69段。两个基座上的答案覆盖率分别为0.770和0.810，高于固定预算，说明它更少弃权。

确定分支因此严格优于固定16段基线：准确率更高、证据更少。规划器单独运行时成本更低，但在GLM上留下小缺口；置信度门控回退在不对每个问题支付16段成本的情况下补上了缺口。

\subsection{与具名RAG基线比较}

表\ref{tab:docfinqa}同时报告四个具名RAG基线。DeepSeek上，Self-RAG与CRAG取得0.373，Adaptive-RAG取得0.363，FLARE取得0.330；GLM上，Self-RAG与CRAG取得0.447，Adaptive-RAG与FLARE取得0.393。FinPlan+fallback在DeepSeek上比最好的具名基线高0.084，在GLM上高0.050，且平均段落数不高于这些基线。四个基线保留各自的检索或反思策略，使用与FinPlan相同的段落块、阅读器和评分，因此差异来自证据策略而不是阅读器。

Qwen3-8B的DocFinQA结果放在附录\ref{app:qwen}：所有条件的精确匹配都低于0.11，主要原因是阅读器返回解释性文字而不是短答案，该基座不用于确定分支的主比较。

\begin{table}[htbp]
\centering
\zihao{5}
\setlength{\tabcolsep}{2pt}
\caption{DocFinQA长文档问答结果。Acc.为精确匹配，Macro F1为宏F1，Paras为平均段落，Cov.为覆盖率。}
\label{tab:docfinqa}
\begin{tabular}{lrrrrrrrr}
\toprule
 & \multicolumn{4}{c}{DeepSeek Flash} & \multicolumn{4}{c}{GLM 5.3 Flash} \\
\cmidrule(lr){2-5}\cmidrule(lr){6-9}
Condition & Acc. & Macro F1 & Paras & Cov. & Acc. & Macro F1 & Paras & Cov. \\
\midrule
Fixed-4 & .333 & .192 & 4.00 & .557 & .400 & .199 & 4.00 & .653 \\
Fixed-8 & .367 & .223 & 8.00 & .603 & .440 & .225 & 8.00 & .720 \\
Fixed-16 & .390 & .234 & 16.00 & .640 & .480 & .249 & 16.00 & .757 \\
Planner B=8 & .417 & .230 & 4.56 & .707 & .450 & .248 & 4.15 & .763 \\
Planner B=16 & .460 & .276 & 6.85 & .773 & .470 & .243 & 5.88 & .840 \\
Self-RAG & .373 & .223 & 8.03 & .613 & .447 & .231 & 8.16 & .717 \\
CRAG & .373 & .222 & 8.05 & .613 & .447 & .235 & 8.32 & .727 \\
FLARE & .330 & .197 & 4.00 & .557 & .393 & .193 & 4.19 & .643 \\
Adaptive-RAG & .363 & .211 & 7.22 & .570 & .393 & .198 & 6.51 & .657 \\
\midrule
\textbf{FinPlanRAG (obligation + fallback)} & \textbf{.457} & \textbf{.276} & 10.69 & \textbf{.770} & \textbf{.497} & \textbf{.252} & 10.31 & \textbf{.810} \\
\bottomrule
\end{tabular}
\end{table}


\section{模糊分支：227条核验断言}

\subsection{实验口径}

模糊分支在227条XBRL核验断言上评估。验证集组合比较校准随机森林后验与“LLM锚点后验 + RF”的融合后验，选择随机森林后验作为最终点预测和决策头；融合权重$w=0.1$的锚点后验在消融中报告。LLM锚点抽取器在三个基座上运行，最终RF决策头不调用LLM，三个基座上的点指标相同。表\ref{tab:vague}报告三个基座上的主结果，完整十一项指标见附录表\ref{tab:vague-ds}至表\ref{tab:vague-qwen}。

\subsection{主结果}

表\ref{tab:vague}按DeepSeek、GLM和Qwen3三个基座分组报告Acc、Bal、Macro和Loss。FinPlanRAG在三个基座上取得0.533准确率、0.491平衡准确率、0.486宏F1和3.454决策损失。FinDVer在DeepSeek上取得0.551、0.461、0.464和4.678，在GLM上取得0.498、0.439、0.442和5.233；Fin-Fact在Qwen3上取得0.511、0.477、0.464和4.846。

普通准确率在GLM和Qwen3上高于最强具名基线，在DeepSeek上与FinDVer统计平局。平衡准确率和宏F1的点估计在三个基座上均更高，但配对区间包含零。决策损失在三个基座上均显著更低，配对区间排除零。校准指标ECE和Brier在三个基座上均优于全部具名基线，说明RF后验的概率质量更可靠。

\begin{table}[htbp]
\centering
\zihao{5}
\setlength{\tabcolsep}{2.5pt}
\caption{227条XBRL核验断言上的主结果。FinPlanRAG为包含LLM锚点信号的AnchorEvolve-Fusion，共享校准RF仅作非LLM诊断。Acc为准确率，Bal为平衡准确率，Macro为宏F1，Loss为非对称决策损失（越低越好）。完整指标见附录表\ref{tab:vague-ds}至表\ref{tab:vague-qwen}。}
\label{tab:vague}
\resizebox{\linewidth}{!}{%
\begin{tabular}{lrrrrrrrrrrrr}
\toprule
方法 & DeepSeek Acc & Bal & Macro & Loss & GLM Acc & Bal & Macro & Loss & Qwen Acc & Bal & Macro & Loss \\
\midrule
FinDVer-style & \textbf{.551} & .461 & .464 & 4.678 & .498 & .439 & .442 & 5.233 & .445 & .422 & .406 & 5.524 \\
Fin-Fact-style & .537 & .446 & .450 & 4.705 & .454 & .398 & .402 & 5.339 & .511 & .477 & .464 & 4.846 \\
Self-RAG-style & .546 & .451 & .456 & 4.564 & .454 & .406 & .404 & 5.163 & .489 & .447 & .437 & 5.093 \\
CRAG-style & .529 & .423 & .418 & 4.802 & .515 & .431 & .435 & 4.846 & .454 & .418 & .399 & 5.366 \\
FLARE-style & .511 & .417 & .411 & 4.881 & .471 & .388 & .387 & 5.286 & .326 & .298 & .212 & 6.740 \\
Adaptive-RAG-style & .493 & .406 & .400 & 5.093 & .427 & .372 & .376 & 5.630 & .344 & .324 & .290 & 6.485 \\
\midrule
AnchorEvolve-LLM & .445 & .307 & .242 & 4.035 & .454 & .329 & .295 & 3.903 & .427 & .353 & .355 & 3.586 \\
\midrule
FinPlanRAG（AnchorEvolve-Fusion, $w=0.1$） & .524 & .471 & .468 & 3.604 & \textbf{.537} & .489 & \textbf{.487} & 3.648 & \textbf{.537} & .486 & .485 & 3.577 \\
共享校准RF（非LLM诊断） & .533 & \textbf{.491} & \textbf{.486} & \textbf{3.454} & .533 & \textbf{.491} & .486 & \textbf{3.454} & .533 & \textbf{.491} & \textbf{.486} & \textbf{3.454} \\
\bottomrule
\end{tabular}
}
\end{table}


\subsection{综合比较}

结果呈现的是语义融合与风险校准的分工，而不是统一的分类收益。FinPlanRAG的融合后验保留不同LLM的锚点判断，并在GLM 5.3 Flash和Qwen3-8B上取得最高普通准确率；共享校准RF作为非LLM诊断，在三套基座上产生相同点指标并取得最低决策损失。LLM锚点层负责提供证据资格和审计轨迹，校准后端负责将概率状态转换为风险可控的输出。

\section{路由与混合负载}

\subsection{路由准确率}

表\ref{tab:router}报告573条混合负载。DeepSeek路由器正确路由570条（0.995）：300条确定问题中297条，273条模糊断言全部正确。三条确定问题被送到模糊分支。GLM路由器正确路由567条（0.990）：297条确定问题、270条模糊断言。Qwen3路由器正确路由271条确定问题和226条模糊断言（0.867），错误方向均衡（29条确定问题送到模糊分支，47条模糊断言送到确定分支），说明8B基座不能稳定判断语义边界。

\subsection{混合效用}

DeepSeek的混合效用为0.403，等于oracle，高于Always-RAG的0.121和Always-Anchor的0.282。GLM的混合效用为0.419，对oracle的0.422、Always-RAG的0.130、Always-Anchor的0.292，恢复99.3\%的oracle效用。Qwen3的混合效用为0.263，低于Always-Anchor的0.298，说明小基座不适合承担路由器角色。路由结果按基座报告，不合并。

\begin{table}[htbp]
\centering
\zihao{5}
\setlength{\tabcolsep}{3pt}
\caption{混合负载路由结果。效用定义见第\ref{ch:experiments}章。}
\label{tab:router}
\resizebox{\linewidth}{!}{%
\begin{tabular}{lrrrrr}
\toprule
基座 & 路由准确率 & 路由器 & 始终RAG & 始终Anchor & Oracle \\
\midrule
DeepSeek & .995 & .403 & .121 & .282 & .403 \\
GLM 5.3 & .990 & .419 & .130 & .292 & .422 \\
Qwen3-8B & .867 & .263 & .018 & .298 & .316 \\
\bottomrule
\end{tabular}
}
\end{table}


\section{消融}

\subsection{融合权重}

表\ref{tab:ablation}变化融合权重。纯LLM锚点后验（$w=1$）在DeepSeek和GLM上的平衡准确率为0.307和0.329，明显偏向高于桶。纯随机森林后验（$w=0$）为0.491和0.491。融合$w=0.1$在三个基座上的平衡准确率为0.471、0.489、0.486，宏F1为0.468、0.487、0.485，均不高于纯随机森林后验；更大的权重进一步降低校准。

\subsection{决策头选择}

验证集组合因此选择随机森林头，其决策损失为3.454；融合头的决策损失为3.604、3.648、3.577。消融说明两个效应：随机森林提供校准和点判别；LLM锚点提供语义资格轨迹和可审计的硬软约束，但其纯后验不足以替代校准后验。最终系统保留锚点抽取用于审计和证据资格，而决策头使用验证集选中的随机森林后验。

\begin{table*}[htbp]
\centering
\zihao{5}
\setlength{\tabcolsep}{3pt}
\caption{锚点后验融合权重$w$的消融。$w=0$为FinPlan-RF，$w=1$为纯LLM锚点后验。}
\label{tab:ablation}
\resizebox{0.96\textwidth}{!}{%
\begin{tabular}{lrrrrrrrr}
\toprule
方法 & DeepSeek Acc & Bal & Macro & Loss & GLM Acc & Bal & Macro & Loss \\
\midrule
AnchorEvolve-LLM ($w=1$) & .445 & .307 & .242 & 4.035 & .454 & .329 & .295 & 3.903 \\
FinPlan-RF ($w=0$) & \textbf{.533} & \textbf{.491} & \textbf{.486} & \textbf{3.454} & .533 & \textbf{.491} & .486 & \textbf{3.454} \\
融合 $w=.1$ & .524 & .471 & .468 & 3.604 & \textbf{.537} & .489 & \textbf{.487} & 3.648 \\
融合 $w=.2$ & .511 & .452 & .453 & 3.639 & .529 & .474 & .477 & 3.683 \\
融合 $w=.3$ & .520 & .449 & .450 & 3.718 & \textbf{.537} & .465 & .473 & 3.674 \\
融合 $w=.4$ & .507 & .406 & .400 & 3.762 & .524 & .425 & .429 & 3.595 \\
\bottomrule
\end{tabular}
}
\end{table*}


\section{外部迁移：FinDVer}

表\ref{tab:findver}报告FinDVer testmini默认阅读器结果。DeepSeek上规划器加回退取得0.701准确率，平均8.76段；固定16段0.660，固定8段0.623。GLM上该分支取得0.753，平均9.19段；固定16段0.743，固定8段0.680。Qwen3上该分支取得0.649，平均6.13段；固定16段0.636，固定8段0.629。分支在三个基座上均高于固定16段阅读器，平均少用6.9至9.9段。官方思维链阅读器在DeepSeek和GLM上分别取得0.817和0.819（金证据），十段检索版本均为0.787；这些使用不同阅读器和证据接口，是上界参考，不是同条件基线。FinDVer是迁移集，分支未在其上调参。

\begin{table}[htbp]
\centering
\zihao{5}
\setlength{\tabcolsep}{3pt}
\caption{FinDVer testmini迁移结果（默认阅读器）。}
\label{tab:findver}
\begin{tabular}{llrrr}
\toprule
基座 & 条件 & 准确率 & 宏F1 & 平均段落 \\
\midrule
\multicolumn{5}{l}{\textit{DeepSeek}} \\
固定4段 & .551 & .665 & 4.00 \\
固定8段 & .623 & .713 & 8.00 \\
固定16段 & .660 & .735 & 16.00 \\
规划器 & .680 & .752 & 5.63 \\
规划器+回退 & .701 & .759 & 8.76 \\
\midrule
\multicolumn{5}{l}{\textit{GLM 5.3}} \\
固定4段 & .599 & .701 & 4.00 \\
固定8段 & .680 & .753 & 8.00 \\
固定16段 & .743 & .792 & 16.00 \\
规划器 & .726 & .774 & 5.69 \\
规划器+回退 & .753 & .789 & 9.19 \\
\midrule
\multicolumn{5}{l}{\textit{Qwen3-8B}} \\
固定4段 & .593 & .635 & 4.00 \\
固定8段 & .629 & .664 & 8.00 \\
固定16段 & .636 & .673 & 16.00 \\
规划器 & .637 & .657 & 5.12 \\
规划器+回退 & .649 & .663 & 6.13 \\
\bottomrule
\end{tabular}
\end{table}


\section{误差分析}

\subsection{确定分支}

确定分支的主要错误是证据缺失和算术错误。DocFinQA的段落块可能把表格行拆开，导致阅读器找不到分子或分母；规划器的\texttt{missing}字段可以部分修复，但当缺失项不在BM25候选的前32条中时无法恢复。Qwen3在DocFinQA上的精确匹配接近地板，主要原因是回答包含解释性文字，而不是数值本身错误。

\subsection{模糊分支}

模糊分支的主要错误是锚点抽取器过度预测高于桶。纯LLM后验在DeepSeek上把大量断言判为高于，导致平衡准确率只有0.307。随机森林通过结构化特征纠正了大部分偏差；最终系统选择RF头后，LLM锚点不再影响点预测。硬锚点很少被抽取器触发，触发时可靠性门槛和冲突检查可以避免错误排除，但也限制了硬约束的覆盖。可信集在大多数断言上包含两个或三个桶，覆盖率接近0.99，集合宽度是下一步需要压缩的对象。

\subsection{路由器}

路由器的错误集中在语义边界。DeepSeek和GLM把“未来支付金额”“预期股利”等确定问题判为模糊断言；这些问题的答案在文件中已经给出，但文本包含未来时态。Qwen3的错误更多，说明小模型的语义边界判断不稳定。提高路由器的一个方向是给路由器提供文件中的候选段落标题，而不是只给查询文本；本文不采用该方向，以保持路由成本为一次短调用。
